\section{Introduction}

Small unmanned aerial vehicles (UAVs) are increasingly used in logistics, inspection, emergency response, and low-altitude sensing. The same accessibility also creates a demand for vision-based detection around protected airspace. Compared with generic object detection, long-range UAV detection is unusually unforgiving: a target may cover fewer than $16\times16$ pixels after letterboxing, motion and defocus weaken its contours, and clouds, birds, buildings, and foliage generate visually similar responses. A detector must therefore raise recall without flooding the scene with false positives, and it must do so under the memory and latency constraints of an edge platform. Recent surveys identify this accuracy--latency tension as a central problem in vision-based anti-UAV systems \cite{yasmeen2025review,wang2021smallobjectsurvey}.

One-stage detectors such as YOLO and SSD established direct dense prediction as a practical route to real-time detection \cite{redmon2016yolo,liu2016ssd}. Feature-pyramid and path-aggregation designs subsequently improved the exchange of high-level semantics and fine spatial details \cite{lin2017fpn,liu2018panet}. Anchor-free formulations and quality-aware classification further reduced hand-designed priors and aligned confidence with localization quality \cite{tian2019fcos,lin2017focal,li2020gfl}. Yet simply increasing the number of high-resolution branches or stacking attention modules can increase memory traffic and kernel-launch overhead. EfficientDet, YOLOv4, and YOLOv7 demonstrate that practical speed depends on coordinated backbone, neck, training, and deployment design rather than FLOPs alone \cite{tan2020efficientdet,bochkovskiy2020yolov4,wang2023yolov7}.

The central observation of this study is that classification and localization do not require identical feature subspaces. Classification favors semantic invariance and inter-class discrimination, whereas localization must preserve boundary- and displacement-sensitive information. TSD showed that a common proposal and representation can induce task conflict \cite{wu2020tsd}, and TOOD demonstrated the value of combining task interaction with task-specific features and alignment \cite{feng2021tood}. A fully decoupled head reduces direct feature competition but duplicates costly spatial processing. Conversely, a fully shared head is compact but constrains both tasks to the same representation budget. These considerations motivate a constrained alternative: retain a compact basis for both tasks and allocate a small additional representation only to classification.

We therefore propose \method{}, short for Shared-Representation and Private-Rank Allocation YOLO, as a structural modification of the YOLOv8 one-stage framework. YOLOv8n serves as the primary baseline for attributing changes in accuracy, parameter count, and latency. The detector retains P2--P4 feature maps through a lightweight bidirectional neck and applies a structurally re-parameterized head at each scale. A 64-channel shared stem supports box-distribution and class prediction, whereas a 16-channel classification-private stem supplies an additive class residual. Zero initialization of the private output projection preserves the shared detector's initial input--output function. Before inference, RepVGG-style structural re-parameterization \cite{ding2021repvgg} concatenates the equivalent stem kernels and merges their task projections into a single path. Classification capacity is therefore expanded during optimization without retaining a parallel private inference branch.

Training follows two stages. Stage I learns the shared detector from standard supervision and response distillation. Distribution-level localization distillation transfers the teacher's discrete box distributions \cite{zheng2022ld}, while a confidence-weighted term emphasizes reliable teacher locations. Stage II starts from the function-preserving zero-residual state and optimizes only the classification-private stem and projection. Under this constraint, the learned box function remains fixed while classification receives additional representational capacity.

The method combines constrained task-specific representation allocation, function-preserving initialization, and exact deployment fusion; task decoupling, zero initialization, distillation, and structural re-parameterization are established techniques individually. The principal contributions are as follows:
\begin{enumerate}
\item We formulate prediction-head capacity as shared and classification-private representation allocation. The resulting 64+16 design makes the box output directly independent of the private representation, while allowing classification to use the combined shared and private row spaces.
\item We introduce a function-preserving two-stage realization in which a zero-output private projection is refined after the shared detector has converged. This isolates additional classification capacity without changing the established box-regression mapping during Stage II.
\item We derive an exact fusion rule that converts the shared and private training paths into one $3\times3$ stem and one $1\times1$ output layer per scale, preserving the computed outputs up to floating-point roundoff and eliminating private-branch inference overhead.
\item A controlled evaluation with validation-selected checkpoints, IoU-resolved AP, matched Stage-II seeds 40/41/42, and same-process paired latency measurement characterizes the resulting detector. Relative to YOLOv8n, the final checkpoint improves Precision, Recall, AP@0.50, \mapall{}, and AP@0.75 by 3.088, 8.242, 5.057, 6.895, and 9.750 percentage points, respectively, while reducing deployed parameters by 28.0\%.
\end{enumerate}
